% Source PDF: source/普通高考/2014/2014新课标2文(贵州,甘肃,青海,西藏,黑龙江,吉林,海南,宁夏,内蒙古,新疆,云南).pdf, page 1, question 8.
% pdfimages -list found no embedded bitmap; the program flowchart is vector line art.
% Grid evidence: tmp/review_flowchart_2014_new_curriculum_paper_2_liberal/grid/page-1-grid100.png
% (100px coarse + 50px fine) and q8-block-grid50.png (50px coarse + 25px fine).
% Comparison crop: tmp/review_flowchart_2014_new_curriculum_paper_2_liberal/crop/q8_original_final.png,
% taken from the ungridded page render with a safety border.
% Elements: rounded start/end terminals; input trapezoid; M=1,S=3 and k=1
% process rectangles; k<=t decision; M=(M/k)x, S=M+S, k=k+1 rectangles;
% output parallelogram; 是/否 labels; and the left loop-back connector.
% Relations: start -> input -> M=1,S=3 -> k=1 -> decision; 是 -> M=(M/k)x
% -> S=M+S -> k=k+1 -> decision; 否 -> 输出 S -> 结束. The initialization
% stem and loop return meet before the decision; the loop's horizontal segment
% points right, followed by one downward arrow into the diamond. No dashed
% segments or shared output-loop connections.
\begin{tikzpicture}[
  x=1cm,y=1cm,baseline,
  flowarrow/.style={-{Latex[length=2.1mm,width=1.5mm]},line width=0.45pt},
  nodebox/.style={draw,line width=0.45pt,inner sep=2pt,minimum height=0.66cm,
    anchor=center,align=center,execute at begin node={\setlength{\parindent}{0pt}}},
  branchlabel/.style={anchor=center,align=center,inner sep=0pt,
    execute at begin node={\setlength{\parindent}{0pt}}},
  term/.style={nodebox,rounded corners=2.5pt,minimum width=1.30cm,text width=1.30cm},
  io/.style={nodebox,shape=trapezium,trapezium left angle=70,trapezium right angle=110,
    minimum width=2.00cm,text width=1.70cm,minimum height=0.72cm},
  decision/.style={nodebox,diamond,aspect=2.7,minimum width=3.75cm,minimum height=1.02cm,inner sep=1pt}
]
  % Keep a narrow, even margin around the left return loop and output branch.
  \path[use as bounding box] (-4.30,0.20) rectangle (5.30,11.45);

  \node[term] (start) at (0,11.10) {开始};
  \node[io] (input) at (0,9.82) {输入 $x,t$};
  \node[nodebox,text width=2.90cm,minimum width=2.90cm] (init) at (0,8.52) {$M=1,S=3$};
  \node[nodebox,text width=1.45cm,minimum width=1.45cm] (kinit) at (0,7.22) {$k=1$};
  \node[decision] (test) at (0,5.72) {$k\leq t$};
  \node[nodebox,text width=2.50cm,minimum width=2.50cm,minimum height=1.05cm] (multiply) at (-2.30,4.05)
    {$M=\dfrac{M}{k}x$};
  \node[nodebox,text width=2.45cm,minimum width=2.45cm] (sum) at (-2.30,2.45) {$S=M+S$};
  \node[nodebox,text width=2.10cm,minimum width=2.10cm] (increment) at (-2.30,0.85) {$k=k+1$};
  \node[io] (output) at (3.35,4.95) {输出 $S$};
  \node[term] (finish) at (3.35,3.35) {结束};

  % Main initialization stem; it joins the loop junction without an extra arrowhead.
  \draw[flowarrow] (start.south) -- (input.north);
  \draw[flowarrow] (input.south) -- (init.north);
  \draw[flowarrow] (init.south) -- (kinit.north);
  \draw (kinit.south) -- (0,6.62);
  \draw[flowarrow] (0,6.62) -- (test.north);

  % Affirmative branch through the three process nodes.
  \draw[flowarrow] (test.west) -- node[branchlabel,above] {是} (-2.30,5.72) -- (multiply.north);
  \draw[flowarrow] (multiply.south) -- (sum.north);
  \draw[flowarrow] (sum.south) -- (increment.north);

  % Negative branch to output and termination, independent of the loop route.
  \draw[flowarrow] (test.east) -- node[branchlabel,above] {否} (3.35,5.72) -- (output.north);
  \draw[flowarrow] (output.south) -- (finish.north);

  % Return path from k=k+1 to the single input point above the decision.
  \draw (increment.west) -- (-4.00,0.85) -- (-4.00,6.62);
  \draw[flowarrow] (-4.00,6.62) -- (-0.60,6.62) -- (0,6.62);

\end{tikzpicture}
